\documentclass{beamer}
\usetheme{Madrid} % Sets basic formatting.
\usecolortheme{dolphin}	% Sets the colour scheme.

%% packages
\usepackage{amsmath}
\usepackage{amssymb}
\usepackage{bm}
\usepackage{booktabs}
\usepackage{tikz}
\usetikzlibrary{positioning,arrows.meta,calc}
\usepackage{textcomp}

\usepackage[style=ieee,backend=biber]{biblatex}
\addbibresource{references.bib}
\renewcommand{\footnotesize}{\scriptsize}

\graphicspath{{figs/}}

\setbeamertemplate{navigation symbols}{} % Manually changes one piece of formatting. 

\title[Practical Radar Target Detection in K-distributed Clutter]{Practical Radar Target Detection in K-distributed Clutter}
\author[]{Tian Huang}
\institute[]{Supervisor: Prof. Elias Aboutanios}
\date{November 9, 2026}

\begin{document}

% Presentations come in slide frames.  You have to tell LaTeX when to start a frame, and when to end the frame.  The most common error beginners make with beamer is forgetting the \end{frame} command.	

%% ----- Title ------ %%
\begin{frame}	
    \titlepage	% Prints a title page populated with the information given in the preamble
\end{frame}		

%% ----- Introduction ------ %%
\begin{frame}{Introduction}
    \begin{itemize}
        \item What is target detection?
        \item In what applications is target detection used?
        \item However, in an example such as a maritime environment, the sea can return more power than the target. This is known as clutter (i.e. define clutter), in addition to noise (define noise).
        \item Problem statement: Use definition
        \item Effective target target detection requires the suppression of clutter and interference. This is achieved by applying adaptive filters to the cell under test (CUT).
    \end{itemize}

    \pause
    
    \begin{block}{The Problem}
        \begin{itemize}
            \item Outline the problem. How can we detect a weak target in strong, spiky, changing clutter that is not Gaussian distributed?
        \end{itemize}     
    \end{block}
\end{frame}

%% ----- Outline ------ %%
\begin{frame}{Outline}
  \tableofcontents  
\end{frame}

%% ----- Section 1: Background ------ %%
\section{Background}

%% ----- Section 1.1: Airborne Radar Signal Model ------ %%
\subsection{Airborne Radar Signal Model}
\begin{frame}{Airborne Radar Signal Model}
    
\end{frame}

%% ----- Section 1.2: Space-Time Adaptive Processing (STAP) ------ %%
\subsection{Space-Time Adaptive Processing (STAP)}
\begin{frame}{Space-Time Adaptive Processing (STAP)}
    \begin{itemize}
        \item Dimension 1:
        \item Dimension 2: Maybe do this in columns \cite{Aboutanios2020}
    \end{itemize}
\end{frame}

%% ----- Section 1.3: K-distributed Clutter Model ------ %%
\subsection{K-distributed Clutter}
\begin{frame}{K-distributed Clutter}
    \begin{itemize}
        \item Together, what our received vector looks like.
        \item Covariance matrices, etc. Explicit that they are not retrievable in practice and are instead estimated.
    \end{itemize}
\end{frame}

%% ----- Section 2: Traditional Adaptive Detectors ------ %%
\section{Traditional Adaptive Detectors}
\begin{frame}{Traditional Adaptive Detectors}
    \begin{itemize}
        \item Detection theory
        \item Traditional detectors: AMF, NAMF,KGLRT
        \item But sea-clutter is not Gaussian so classical detectors struggle.
    \end{itemize}
\end{frame}

%% ----- Section 3: Current Work ------ %%
\section{K-distributed Clutter Detectors}
\begin{frame}{K-distributed Clutter Detectors}
    \begin{itemize}
        \item Gini's OKD
        \item TEN detector
    \end{itemize}
\end{frame}

%% ----- Section 4: Research Gap and Aim ------ %%
\section{Research Gap and Aim}
\begin{frame}{Research Gap and Aim}
    \begin{itemize}
        \item Gap
        \item Observations
        \item Aims
        \item Approach
    \end{itemize}
\end{frame}

%% ----- Section 5: Preliminary Work ------ %%
\section{Preliminary Work}
\begin{frame}{Preliminary Work}
    \begin{itemize}
        \item Preliminary work
        \item Gantt chart (also set 3 goals - one for each term)
        \item Risk assessment, skills to be learnt, goals, objectives and tasks done to achieve goals.
    \end{itemize}
\end{frame}

%% ----- Section 4: References ------ %%
\section{References}
\begin{frame}{References}
    \printbibliography[heading=none]
\end{frame}




% \begin{frame}{Uniform convergence}	

% \begin{definition}	% There are lots of "theorem-like" environments for beamer just as 	usual with LaTeX: definition, theorem, lemma, example, proof, etc...

% A sequence of functions $f_n \colon \mathbb{R} \to \mathbb{R}$ \emph{converges uniformly} to a function $f \colon \mathbb{R} \to \mathbb{R}$ if for all $\epsilon > 0$ there exists an $N \in \mathbb{N}$ such that $n \geq N$ implies 
% \[
% \sup_{x \in \mathbb{R}} |f_n(x) - f(x)| < \epsilon.
% \]

% \end{definition}

% \pause	% Generates a break in the slide presentation

% \begin{block}{Pointwise and uniform continuity} % Blocks are a beamer speciality.

% \begin{itemize}
% \item Uniform convergence implies pointwise convergence
% \item Pointwise convergence does not imply uniform convergence
% \end{itemize}

% \end{block}

% \pause

% \begin{theorem}
% Let $f_n \colon\mathbb{R} \to \mathbb{R}$ be continuous and converge uniformly to $f \colon \mathbb{R} \to \mathbb{R}$.  Then $f$ is continuous.
% \end{theorem}

% \end{frame}

% \begin{frame}{Uniform convergence and continuity}

% \begin{proof}
% Let $x \in \mathbb{R}$ and let $\epsilon > 0$.  There exists $N \in \mathbb{N}$ such that $n \geq N$ implies
% \begin{equation}
% \label{eq1}
% \sup_{x \in \mathbb{R}} |f_n(x) - f(x)| < \frac{\epsilon}{3}.
% \end{equation}

% There exists $\delta > 0$ such that 
% \begin{equation}
% \label{eq2}
% |f_N(x) - f_N(y) | < \frac{\epsilon}{3}\ \mbox{whenever}\ |x-y| < \delta.
% \end{equation}

% Then inequalities \eqref{eq1} and \eqref{eq2} imply that whenever $|x-y | < \delta$, we have 
% \begin{align*}
% |f(x) - f(y)| 
% & \leq |f(x) - f_N(x)| + |f_N(x) - f_N(y)| + | f_N(y) - f(y) | \\
% & < \frac{\epsilon}{3} + \frac{\epsilon}{3} + \frac{\epsilon}{3} \\
% & = \epsilon.
% \end{align*} 
% \end{proof}

% \end{frame}

% \begin{frame}{People}

% This subject owes much to 

% \begin{columns}	% A handy way of putting things side-by-side

% \begin{column}{0.4\textwidth} % this says the column will be the width of 0.4 times the \textwidth
% \begin{figure}	% for graphics
% \includegraphics[height=4cm]{Cauchy}	% save your image file in your project, and include using \includegraphics{filename}.  Specify height, width, scaling as shown.
% \caption{Augustin-Louis Cauchy}	% there are ways to delete the "Figure:" if you want to
% \end{figure}
% \end{column}

% \begin{column}{0.4\textwidth}
% \begin{figure}
% \includegraphics[height=4cm]{Weierstrass}
% \caption{Karl Weierstrass}
% \end{figure}
% \end{column}

% \end{columns}

% \end{frame}

\end{document}	